\documentclass{amsart}
% For dealing with references we use the comment environment
\usepackage{verbatim}
\newenvironment{reference}{\comment}{\endcomment}
%\newenvironment{reference}{}{}
\newenvironment{slogan}{\comment}{\endcomment}
\newenvironment{history}{\comment}{\endcomment}

% For commutative diagrams we use Xy-pic
\usepackage[all]{xy}

% We use 2cell for 2-commutative diagrams.
\xyoption{2cell}
\UseAllTwocells

% We use multicol for the list of chapters between chapters
\usepackage{multicol}

% This is generally recommended for better output
\usepackage{lmodern}
\usepackage[T1]{fontenc}

% For cross-file-references

% Package for hypertext links:
\usepackage{hyperref}

% For any local file, say "hello.tex" you want to link to please

% Theorem environments.
%
\theoremstyle{plain}
\newtheorem{theorem}[subsection]{Theorem}
\newtheorem{proposition}[subsection]{Proposition}
\newtheorem{lemma}[subsection]{Lemma}

\theoremstyle{definition}
\newtheorem{definition}[subsection]{Definition}
\newtheorem{example}[subsection]{Example}
\newtheorem{exercise}[subsection]{Exercise}
\newtheorem{situation}[subsection]{Situation}

\theoremstyle{remark}
\newtheorem{remark}[subsection]{Remark}
\newtheorem{remarks}[subsection]{Remarks}

\numberwithin{equation}{subsection}

% Macros
%
\def\lim{\mathop{\mathrm{lim}}\nolimits}
\def\colim{\mathop{\mathrm{colim}}\nolimits}
\def\Spec{\mathop{\mathrm{Spec}}}
\def\Hom{\mathop{\mathrm{Hom}}\nolimits}
\def\Ext{\mathop{\mathrm{Ext}}\nolimits}
\def\SheafHom{\mathop{\mathcal{H}\!\mathit{om}}\nolimits}
\def\SheafExt{\mathop{\mathcal{E}\!\mathit{xt}}\nolimits}
\def\Sch{\mathit{Sch}}
\def\Mor{\mathop{\mathrm{Mor}}\nolimits}
\def\Ob{\mathop{\mathrm{Ob}}\nolimits}
\def\Sh{\mathop{\mathit{Sh}}\nolimits}
\def\NL{\mathop{N\!L}\nolimits}
\def\CH{\mathop{\mathrm{CH}}\nolimits}
\def\proetale{{pro\text{-}\acute{e}tale}}
\def\etale{{\acute{e}tale}}
\def\QCoh{\mathit{QCoh}}
\def\Ker{\mathop{\mathrm{Ker}}}
\def\Im{\mathop{\mathrm{Im}}}
\def\Coker{\mathop{\mathrm{Coker}}}
\def\Coim{\mathop{\mathrm{Coim}}}

% Boxtimes
%
\DeclareMathSymbol{\boxtimes}{\mathbin}{AMSa}{"02}

%
% Macros for moduli stacks/spaces
%
\def\QCohstack{\mathcal{QC}\!\mathit{oh}}
\def\Cohstack{\mathcal{C}\!\mathit{oh}}
\def\Spacesstack{\mathcal{S}\!\mathit{paces}}
\def\Quotfunctor{\mathrm{Quot}}
\def\Hilbfunctor{\mathrm{Hilb}}
\def\Curvesstack{\mathcal{C}\!\mathit{urves}}
\def\Polarizedstack{\mathcal{P}\!\mathit{olarized}}
\def\Complexesstack{\mathcal{C}\!\mathit{omplexes}}
% \Pic is the operator that assigns to X its picard group, usage \Pic(X)
% \Picardstack_{X/B} denotes the Picard stack of X over B
% \Picardfunctor_{X/B} denotes the Picard functor of X over B
\def\Pic{\mathop{\mathrm{Pic}}\nolimits}
\def\Picardstack{\mathcal{P}\!\mathit{ic}}
\def\Picardfunctor{\mathrm{Pic}}
\def\Deformationcategory{\mathcal{D}\!\mathit{ef}}

\setlength{\emergencystretch}{3em}
\begin{document}
\title[Stacks Project, Tag 07RN]{1666 --- Noetherian approximation of a quasi-compact quasi-separated scheme extending a prescribed affine-transition finite-type integer model system on a quasi-compact open}
\maketitle
\noindent Copyright (C) 2005--2025 Johan de Jong. Stacks Project authors. Source: \href{https://stacks.math.columbia.edu/tag/07RN}{Tag 07RN}.

\noindent Editorial difficulty note (not author text). Difficulty H3: The theorem extends a prescribed finite-type integer model of a quasi-compact open across an arbitrary qcqs scheme. Its proof builds fibre-product rings, recovers standard affine opens after finite-stage descent, manufactures finite-type subalgebra models, and synchronizes two-level directed gluing with affine transitions. This global relative Noetherian approximation construction is substantially beyond ordinary affine finite-presentation descent..

\noindent Editorial provenance note (not author text). History: excerpt from revision \texttt{a0444\allowbreak 6e57e\allowbreak c1fbc\allowbreak 252a8\allowbreak 71afc\allowbreak ec775\allowbreak 2fb28\allowbreak 07b14}, limits.tex, lines 1087--1252. Reference presentation was adapted; original-author context and core support blocks were retained or added. The mathematical wording is unchanged. Licensed under GNU FDL 1.2 or later; no invariant sections or cover texts. See ../licenses/COPYING. Context blocks reproduce author definitions and supporting results; they are not additional counted problems.

\begin{situation}
\label{situation-descent}
Let $S = \lim_{i \in I} S_i$ be the limit of a directed system of schemes
with affine transition morphisms $f_{i'i} : S_{i'} \to S_i$
(Lemma \ref{lemma-directed-inverse-system-has-limit}).
We assume that $S_i$ is quasi-compact and quasi-separated for all $i \in I$.
We denote $f_i : S \to S_i$ the projection.
We also choose an element $0 \in I$.
\end{situation}

\begin{definition}
\label{properties-definition-ample}
\begin{reference}
\href{https://stacks.math.columbia.edu/bibliography}{Bibliography: EGA (II Definition 4.5.3)}
\end{reference}
Let $X$ be a scheme.
Let $\mathcal{L}$ be an invertible $\mathcal{O}_X$-module.
We say $\mathcal{L}$ is {\it ample} if
\begin{enumerate}
\item $X$ is quasi-compact, and
\item for every $x \in X$ there exists an $n \geq 1$
and $s \in \Gamma(X, \mathcal{L}^{\otimes n})$ such
that $x \in X_s$ and $X_s$ is affine.
\end{enumerate}
\end{definition}

\begin{lemma}
\label{lemma-descend-opens}
In Situation \ref{situation-descent} we have the following:
\begin{enumerate}
\item Given any quasi-compact open $V \subset S = \lim_i S_i$
there exists an $i \in I$ and a quasi-compact open $V_i \subset S_i$
such that $f_i^{-1}(V_i) = V$.
\item Given $V_i \subset S_i$ and $V_{i'} \subset S_{i'}$
quasi-compact opens such that $f_i^{-1}(V_i) = f_{i'}^{-1}(V_{i'})$
there exists an index $i'' \geq i, i'$ such that
$f_{i''i}^{-1}(V_i) = f_{i''i'}^{-1}(V_{i'})$.
\item If $V_{1, i}, \ldots, V_{n, i} \subset S_i$ are quasi-compact
opens and $S = f_i^{-1}(V_{1, i}) \cup \ldots \cup f_i^{-1}(V_{n, i})$
then $S_{i'} = f_{i'i}^{-1}(V_{1, i}) \cup \ldots \cup f_{i'i}^{-1}(V_{n, i})$
for some $i' \geq i$.
\end{enumerate}
\end{lemma}

\begin{proof}
Choose $i_0 \in I$. Note that $I$ is nonempty as the limit is directed.
For convenience we write $S_0 = S_{i_0}$ and $i_0 = 0$.
Choose an affine open covering $S_0 = U_{1, 0} \cup \ldots \cup U_{m, 0}$.
Denote $U_{j, i} \subset S_i$ the inverse image of $U_{j, 0}$
under the transition morphism for $i \geq 0$.
Denote $U_j$ the inverse image of $U_{j, 0}$ in $S$.
Note that $U_j = \lim_i U_{j, i}$ is a limit of affine
schemes.

\medskip\noindent
We first prove the uniqueness statement: Let
$V_i \subset S_i$ and $V_{i'} \subset S_{i'}$
quasi-compact opens such that $f_i^{-1}(V_i) = f_{i'}^{-1}(V_{i'})$.
It suffices to show that $f_{i''i}^{-1}(V_i \cap U_{j, i''})$ and
$f_{i''i'}^{-1}(V_{i'} \cap U_{j, i''})$ become equal
for $i''$ large enough. Hence we reduce to the case
of a limit of affine schemes. In this case write
$S = \Spec(R)$ and $S_i = \Spec(R_i)$ for all $i \in I$.
We may write $V_i = S_i \setminus V(h_1, \ldots, h_m)$
and $V_{i'} = S_{i'} \setminus V(g_1, \ldots, g_n)$.
The assumption means that the ideals
$\sum g_jR$ and $\sum h_jR$ have the same radical
in $R$. This means that $g_j^N = \sum a_{jj'}h_{j'}$ and
$h_j^N = \sum b_{jj'} g_{j'}$ for some $N \gg 0$ and $a_{jj'}$
and $b_{jj'}$ in $R$.
Since $R = \colim_i R_i$ we can chose an index
$i'' \geq i$ such that the equations
$g_j^N = \sum a_{jj'}h_{j'}$ and
$h_j^N = \sum b_{jj'} g_{j'}$ hold in $R_{i''}$ for some
$a_{jj'}$ and $b_{jj'}$ in $R_{i''}$. This implies that
the ideals $\sum g_jR_{i''}$ and $\sum h_jR_{i''}$ have the same radical
in $R_{i''}$ as desired.

\medskip\noindent
We prove existence: If $S_0$ is affine, then $S_i = \Spec(R_i)$ for all
$i \geq 0$ and $S = \Spec(R)$ with $R = \colim R_i$. Then
$V = S \setminus V(g_1, \ldots, g_n)$ for some $g_1, \ldots, g_n \in R$.
Choose any $i$ large enough so that each of the $g_j$ comes from an
element $g_{j, i} \in R_i$ and take
$V_i = S_i \setminus V(g_{1, i}, \ldots, g_{n, i})$.
If $S_0$ is general, then the opens $V \cap U_j$
are quasi-compact because $S$ is quasi-separated. Hence by the
affine case we see that for each $j = 1, \ldots, m$
there exists an $i_j \in I$ and a quasi-compact open
$V_{i_j} \subset U_{j, i_j}$ whose inverse image in $U_j$
is $V \cap U_j$. Set $i = \max(i_1, \ldots, i_m)$
and let $V_i = \bigcup f_{ii_j}^{-1}(V_{i_j})$.

\medskip\noindent
The statement on coverings follows from the uniqueness statement
for the opens $V_{1, i} \cup \ldots \cup V_{n, i}$ and $S_i$ of $S_i$.
\end{proof}


\begin{lemma}
\label{lemma-limit-quasi-affine}
In Situation \ref{situation-descent} if $S$ is quasi-affine, then
for some $i_0 \in I$ the schemes $S_i$ for $i \geq i_0$ are quasi-affine.
\end{lemma}

\begin{proof}
Choose $i_0 \in I$. Note that $I$ is nonempty as the limit is directed.
For convenience we write $S_0 = S_{i_0}$ and $i_0 = 0$.
Let $s \in S$. We may choose an affine open
$U_0 \subset S_0$ containing $f_0(s)$. Since $S$ is quasi-affine
we may choose an element $a \in \Gamma(S, \mathcal{O}_S)$ such
that $s \in D(a) \subset f_0^{-1}(U_0)$, and such that
$D(a)$ is affine. By Lemma \ref{lemma-descend-section}
there exists an $i \geq 0$ such that $a$
comes from an element $a_i \in \Gamma(S_i, \mathcal{O}_{S_i})$.
For any index $j \geq i$ we denote $a_j$
the image of $a_i$ in the global sections of the
structure sheaf of $S_j$.
Consider the opens $D(a_j) \subset S_j$
and $U_j = f_{j0}^{-1}(U_0)$. Note that
$U_j$ is affine and $D(a_j)$ is a quasi-compact open of $S_j$,
see Properties, Lemma \href{https://stacks.math.columbia.edu/tag/01PV}{lemma affine cap s open (Tag 01PV)}
for example. Hence we may apply Lemma \ref{lemma-descend-opens} to the opens
$U_j$ and $U_j \cup D(a_j)$ to conclude that
$D(a_j) \subset U_j$ for some  $j \geq i$.
For such an index $j$ we see that $D(a_j) \subset S_j$ is an affine open
(because $D(a_j)$ is a standard affine open of the affine open $U_j$)
containing the image $f_j(s)$.

\medskip\noindent
We conclude that for every $s \in S$ there exist
an index $i \in I$, and a global section
$a \in \Gamma(S_i, \mathcal{O}_{S_i})$
such that $D(a) \subset S_i$ is an affine open
containing $f_i(s)$. Because $S$ is quasi-compact we
may choose a single index $i \in I$ and global sections
$a_1, \ldots, a_m \in \Gamma(S_i, \mathcal{O}_{S_i})$
such that each $D(a_j) \subset S_i$ is affine open
and such that $f_i : S \to S_i$ has image contained
in the union $W_i = \bigcup_{j = 1, \ldots, m} D(a_j)$.
For $i' \geq i$ set $W_{i'} = f_{i'i}^{-1}(W_i)$.
Since $f_i^{-1}(W_i)$ is all of $S$ we see
(by Lemma \ref{lemma-descend-opens} again)
that for a suitable $i' \geq i$ we
have $S_{i'} = W_{i'}$. Thus we may replace $i$ by
$i'$ and assume that $S_i = \bigcup_{j = 1, \ldots, m} D(a_j)$.
This implies that $\mathcal{O}_{S_i}$ is an ample invertible
sheaf on $S_i$ (see Properties, Definition \ref{properties-definition-ample})
and hence that $S_i$ is quasi-affine, see
Properties, Lemma \href{https://stacks.math.columbia.edu/tag/01QE}{lemma quasi affine O ample (Tag 01QE)}.
Hence we win.
\end{proof}


\begin{lemma}
\label{lemma-descend-section}
In Situation \ref{situation-descent}.
Suppose that $\mathcal{F}_0$ is a quasi-coherent sheaf on $S_0$.
Set $\mathcal{F}_i = f_{i0}^*\mathcal{F}_0$ for $i \geq 0$ and set
$\mathcal{F} = f_0^*\mathcal{F}_0$.
Then
$$
\Gamma(S, \mathcal{F}) = \colim_{i \geq 0} \Gamma(S_i, \mathcal{F}_i)
$$
\end{lemma}

\begin{proof}
Write $\mathcal{A}_j = f_{i0, *} \mathcal{O}_{S_i}$.
This is a quasi-coherent sheaf of $\mathcal{O}_{S_0}$-algebras
(see Morphisms, Lemma \href{https://stacks.math.columbia.edu/tag/01SA}{lemma affine equivalence algebras (Tag 01SA)})
and $S_i$ is the relative spectrum of $\mathcal{A}_i$ over $S_0$.
In the proof of Lemma \ref{lemma-directed-inverse-system-has-limit}
we constructed $S$ as the relative spectrum of
$\mathcal{A} = \colim_{i \geq 0} \mathcal{A}_i$
over $S_0$. Set
$$
\mathcal{M}_i = \mathcal{F}_0 \otimes_{\mathcal{O}_{S_0}} \mathcal{A}_i
$$
and
$$
\mathcal{M} = \mathcal{F}_0 \otimes_{\mathcal{O}_{S_0}} \mathcal{A}.
$$
Then we have $f_{i0, *} \mathcal{F}_i = \mathcal{M}_i$
and $f_{0, *}\mathcal{F} = \mathcal{M}$. Since $\mathcal{A}$
is the colimit of the sheaves $\mathcal{A}_i$ and since tensor
product commutes with directed colimits, we conclude that
$\mathcal{M} = \colim_{i \geq 0} \mathcal{M}_i$.
Since $S_0$ is quasi-compact and quasi-separated we see that
\begin{eqnarray*}
\Gamma(S, \mathcal{F})
& = &
\Gamma(S_0, \mathcal{M}) \\
& = &
\Gamma(S_0, \colim_{i \geq 0} \mathcal{M}_i) \\
& = &
\colim_{i \geq 0} \Gamma(S_0, \mathcal{M}_i) \\
& = &
\colim_{i \geq 0} \Gamma(S_i, \mathcal{F}_i)
\end{eqnarray*}
see Sheaves, Lemma \href{https://stacks.math.columbia.edu/tag/009F}{lemma directed colimits sections (Tag 009F)} and
Topology, Lemma \href{https://stacks.math.columbia.edu/tag/0069}{lemma topology quasi separated scheme (Tag 0069)}
for the middle equality.
\end{proof}


\begin{lemma}
\label{lemma-quasi-affine-finite-type-over-Z}
Let $W$ be a quasi-affine scheme of finite type over
$\mathbf{Z}$. Suppose $W \to \Spec(R)$ is an
open immersion into an affine scheme. There exists a
finite type $\mathbf{Z}$-algebra $A \subset R$
which induces an open immersion $W \to \Spec(A)$.
Moreover, $R$ is the directed colimit of such subalgebras.
\end{lemma}

\begin{proof}
Choose an affine open covering $W = \bigcup_{i = 1, \ldots, n} W_i$
such that each $W_i$ is a standard affine open in $\Spec(R)$.
In other words, if we write $W_i = \Spec(R_i)$
then $R_i = R_{f_i}$ for some $f_i \in R$.
Choose finitely many $x_{ij} \in R_i$ which generate
$R_i$ over $\mathbf{Z}$.
Pick an $N \gg 0$ such that each $f_i^Nx_{ij}$ comes from an
element of $R$, say $y_{ij} \in R$.
Set $A$ equal to the $\mathbf{Z}$-algebra generated by
the $f_i$ and the $y_{ij}$ and (optionally) finitely many
additional elements of $R$. Then $A$ works. Details omitted.
\end{proof}


\begin{lemma}
\label{lemma-diagram}
Suppose given a cartesian diagram of rings
$$
\xymatrix{
B \ar[r]_s & R \\
B'\ar[u] \ar[r] & R' \ar[u]_t
}
$$
Let $W' \subset \Spec(R')$ be an open of
the form $W' = D(f_1) \cup \ldots \cup D(f_n)$
such that $t(f_i) = s(g_i)$ for some $g_i \in B$
and $B_{g_i} \cong R_{s(g_i)}$. Then $B' \to R'$
induces an open immersion of $W'$ into $\Spec(B')$.
\end{lemma}

\begin{proof}
Set $h_i = (g_i, f_i) \in B'$. More on Algebra,
Lemma \ref{more-algebra-lemma-diagram-localize} shows that
$(B')_{h_i} \cong (R')_{f_i}$ as desired.
\end{proof}


\begin{lemma}
\label{more-algebra-lemma-diagram-localize}
Suppose given a cartesian diagram of rings
$$
\xymatrix{
R &
R' \ar[l]^t \\
B \ar[u]_s &
B'\ar[u] \ar[l]
}
$$
i.e., $B' = B \times_R R'$. If $h \in B'$ corresponds to $g \in B$
and $f \in R'$ such that $s(g) = t(f)$, then the diagram
$$
\xymatrix{
R_{s(g)} = R_{t(f)} &
(R')_f \ar[l]^-t \\
B_g \ar[u]_s &
(B')_h \ar[u] \ar[l]
}
$$
is cartesian too.
\end{lemma}

\begin{proof}
The equality $B' = B \times_R R'$ tells us that
$$
0 \to B' \to B \oplus R' \xrightarrow{s, -t} R
$$
is an exact sequence of $B'$-modules. We have $B_g = B_h$,
$R'_f = R'_h$, and $R_{s(g)} = R_{t(f)} = R_h$ as $B'$-modules.
By exactness of localization
(Algebra, Proposition \href{https://stacks.math.columbia.edu/tag/00CS}{proposition localization exact (Tag 00CS)})
we find that
$$
0 \to B'_h \to B_g \oplus R'_f \xrightarrow{s, -t} R_{s(g)} = R_{t(f)}
$$
is an exact sequence. This proves the lemma.
\end{proof}


\begin{lemma}
\label{lemma-directed-inverse-system-has-limit}
Let $I$ be a directed set. Let $(S_i, f_{ii'})$ be an
inverse system of schemes over $I$. If all the morphisms
$f_{ii'} : S_i \to S_{i'}$ are affine, then the limit $S = \lim_i S_i$ exists
in the category of schemes. Moreover,
\begin{enumerate}
\item each of the morphisms $f_i : S \to S_i$ is affine,
\item for an element $0 \in I$ and any open subscheme $U_0 \subset S_0$
we have
$$
f_0^{-1}(U_0) = \lim_{i \geq 0} f_{i0}^{-1}(U_0)
$$
in the category of schemes.
\end{enumerate}
\end{lemma}

\begin{proof}
Choose an element $0 \in I$. Note that $I$ is nonempty as the limit is
directed. For every $i \geq 0$ consider the quasi-coherent sheaf of
$\mathcal{O}_{S_0}$-algebras $\mathcal{A}_i = f_{i0, *}\mathcal{O}_{S_i}$.
Recall that $S_i = \underline{\Spec}_{S_0}(\mathcal{A}_i)$,
see Morphisms, Lemma \href{https://stacks.math.columbia.edu/tag/01S8}{lemma characterize affine (Tag 01S8)}.
Set $\mathcal{A} = \colim_{i \geq 0} \mathcal{A}_i$.
This is a quasi-coherent sheaf of $\mathcal{O}_{S_0}$-algebras,
see Schemes, Section \href{https://stacks.math.columbia.edu/tag/01LA}{section quasi coherent (Tag 01LA)}.
Set $S = \underline{\Spec}_{S_0}(\mathcal{A})$.
By Morphisms, Lemma \href{https://stacks.math.columbia.edu/tag/01SA}{lemma affine equivalence algebras (Tag 01SA)}
we get for $i \geq 0$ morphisms $f_i : S \to S_i$ compatible with
the transition morphisms. Note that the morphisms $f_i$ are
affine by Morphisms, Lemma \href{https://stacks.math.columbia.edu/tag/01SG}{lemma affine permanence (Tag 01SG)} for example.
By Lemma \href{https://stacks.math.columbia.edu/tag/01YW}{lemma directed inverse system affine schemes has limit (Tag 01YW)} above
we see that for any affine open $U_0 \subset S_0$ the
inverse image $U = f_0^{-1}(U_0) \subset S$ is the limit of the
system of opens $U_i = f_{i0}^{-1}(U_0)$, $i \geq 0$ in the
category of schemes.

\medskip\noindent
Let $T$ be a scheme. Let $g_i : T \to S_i$ be a compatible system
of morphisms. To show that $S = \lim_i S_i$ we have
to prove there is a unique morphism $g : T \to S$ with
$g_i = f_i \circ g$ for all $i \in I$.
For every $t \in T$ there exists an affine open
$U_0 \subset S_0$ containing $g_0(t)$. Let $V \subset g_0^{-1}(U_0)$
be an affine open neighbourhood containing $t$.
By the remarks above we obtain a unique morphism
$g_V : V \to U = f_0^{-1}(U_0)$ such that $f_i \circ g_V = g_i|_{U_i}$
for all $i$. The open sets $V \subset T$ so constructed form
a basis for the topology of $T$. The morphisms $g_V$ glue to a morphism
$g : T \to S$ because of the uniqueness property. This gives the
desired morphism $g : T \to S$.

\medskip\noindent
The final statement is clear from the construction of the limit above.
\end{proof}


\begin{lemma}
\label{lemma-approximate}
Let $S$ be a quasi-compact and quasi-separated scheme. Let $V \subset S$
be a quasi-compact open. Let $I$ be a directed set
and let $(V_i, f_{ii'})$ be an inverse system of schemes over $I$
with affine transition maps, with each $V_i$ of finite type over $\mathbf{Z}$,
and with $V = \lim V_i$. Then there exist
\begin{enumerate}
\item a directed set $J$,
\item an inverse system of schemes $(S_j, g_{jj'})$ over $J$,
\item an order preserving map $\alpha : J \to I$,
\item open subschemes $V'_j \subset S_j$, and
\item isomorphisms $V'_j \to V_{\alpha(j)}$
\end{enumerate}
such that
\begin{enumerate}
\item the transition morphisms $g_{jj'} : S_j \to S_{j'}$ are affine,
\item each $S_j$ is of finite type over $\mathbf{Z}$,
\item $g_{jj'}^{-1}(V'_{j'}) = V'_j$,
\item $S = \lim S_j$ and $V = \lim V'_j$, and
\item the diagrams
$$
\vcenter{
\xymatrix{
V \ar[d] \ar[rd] \\
V'_j \ar[r] & V_{\alpha(j)}
}
}
\quad\text{and}\quad
\vcenter{
\xymatrix{
V'_j \ar[r] \ar[d] & V_{\alpha(j)} \ar[d] \\
V'_{j'} \ar[r] & V_{\alpha(j')}
}
}
$$
are commutative.
\end{enumerate}
\end{lemma}

\begin{proof}
Set $Z = S \setminus V$. Choose affine opens $U_1, \ldots, U_m \subset S$
such that $Z \subset \bigcup_{l = 1, \ldots, m} U_l$. Consider the opens
$$
V \subset V \cup U_1 \subset V \cup U_1 \cup U_2 \subset
\ldots \subset V \cup \bigcup\nolimits_{l = 1, \ldots, m} U_l = S
$$
If we can prove the lemma successively for each of the cases
$$
V \cup U_1 \cup \ldots \cup U_l
\subset
V \cup U_1 \cup \ldots \cup U_{l + 1}
$$
then the lemma will follow for $V \subset S$. In each case we are adding
one affine open. Thus we may assume
\begin{enumerate}
\item $S = U \cup V$,
\item $U$ affine open in $S$,
\item $V$ quasi-compact open in $S$, and
\item $V = \lim_i V_i$ with $(V_i, f_{ii'})$
an inverse system over a directed set $I$, each $f_{ii'}$
affine and each $V_i$ of finite type over $\mathbf{Z}$.
\end{enumerate}
Denote $f_i : V \to V_i$ the projections.
Set $W = U \cap V$. As $S$ is quasi-separated, this is a quasi-compact open
of $V$. By Lemma \ref{lemma-descend-opens}
(and after shrinking $I$) we may assume that there exist
opens $W_i \subset V_i$ such that $f_{ii'}^{-1}(W_{i'}) = W_i$
and such that $f_i^{-1}(W_i) = W$. Since $W$ is a quasi-compact open
of $U$ it is quasi-affine. Hence we may assume (after shrinking $I$ again)
that $W_i$ is quasi-affine for all $i$, see
Lemma \ref{lemma-limit-quasi-affine}.

\medskip\noindent
Write $U = \Spec(B)$. Set $R = \Gamma(W, \mathcal{O}_W)$,
and $R_i = \Gamma(W_i, \mathcal{O}_{W_i})$.
By Lemma \ref{lemma-descend-section} we have $R = \colim_i R_i$.
Now we have the maps of rings
$$
\xymatrix{
B \ar[r]_s & R \\
& R_i \ar[u]_{t_i}
}
$$
We set $B_i = \{(b, r) \in B \times R_i \mid s(b) = t_i(r)\}$ so that we
have a cartesian diagram
$$
\xymatrix{
B \ar[r]_s & R \\
B_i \ar[u] \ar[r] & R_i \ar[u]_{t_i}
}
$$
for each $i$. The transition maps $R_i \to R_{i'}$ induce maps
$B_i \to B_{i'}$. It is clear that $B = \colim_i B_i$.
In the next paragraph we show that for all sufficiently large $i$
the composition $W_i \to \Spec(R_i) \to \Spec(B_i)$ is an open immersion.

\medskip\noindent
As $W$ is a quasi-compact open of $U = \Spec(B)$
we can find a finitely many elements $g_l \in B$, $l = 1, \ldots, m$
such that $D(g_l) \subset W$ and such that
$W = \bigcup_{l = 1, \ldots, m} D(g_l)$.
Note that this implies $D(g_l) = W_{s(g_l)}$ as open subsets of $U$,
where $W_{s(g_l)}$ denotes the largest open subset of $W$ on which
$s(g_l)$ is invertible. Hence
$$
B_{g_l} =
\Gamma(D(g_l), \mathcal{O}_U) =
\Gamma(W_{s(g_l)}, \mathcal{O}_W) = R_{s(g_l)},
$$
where the last equality is
Properties, Lemma \href{https://stacks.math.columbia.edu/tag/01P7}{lemma invert f sections (Tag 01P7)}.
Since $W_{s(g_l)}$ is affine this also
implies that $D(s(g_l)) = W_{s(g_l)}$ as open subsets of $\Spec(R)$.
Since $R = \colim_i R_i$ we can (after shrinking $I$)
assume there exist $g_{l, i} \in R_i$ for all $i \in I$ such that
$s(g_l) = t_i(g_{l, i})$. Of course we choose the $g_{l, i}$
such that $g_{l, i}$ maps to $g_{l, i'}$ under the transition maps
$R_i \to R_{i'}$. Then, by Lemma \ref{lemma-descend-opens} we can
(after shrinking $I$ again)
assume the corresponding opens $D(g_{l, i}) \subset \Spec(R_i)$
are contained in $W_i$ for $l = 1, \ldots, m$ and cover $W_i$.
We conclude that the morphism $W_i \to \Spec(R_i) \to \Spec(B_i)$
is an open immersion, see Lemma \ref{lemma-diagram}.

\medskip\noindent
By Lemma \ref{lemma-quasi-affine-finite-type-over-Z}
we can write $B_i$ as a directed colimit of subalgebras
$A_{i, p} \subset B_i$, $p \in P_i$ each
of finite type over $\mathbf{Z}$ and such that $W_i$ is
identified with an open subscheme of $\Spec(A_{i, p})$.
Let $S_{i, p}$ be the scheme obtained by glueing
$V_i$ and $\Spec(A_{i, p})$ along the open $W_i$, see
Schemes, Section \href{https://stacks.math.columbia.edu/tag/01JA}{section glueing schemes (Tag 01JA)}.
Here is the resulting commutative diagram of schemes:
$$
\xymatrix{
& & V \ar[lld] \ar[d] & W \ar[l] \ar[lld] \ar[d] \\
V_i \ar[d] & W_i \ar[l] \ar[d] & S \ar[lld] & U \ar[lld] \ar[l] \\
S_{i, p} & \Spec(A_{i, p}) \ar[l]
}
$$
The morphism $S \to S_{i, p}$ arises because the upper right
square is a pushout in the category of schemes.
Note that $S_{i, p}$ is of finite type over $\mathbf{Z}$ since
it has a finite affine open covering whose members are
spectra of finite type $\mathbf{Z}$-algebras.
We define a preorder on $J = \coprod_{i \in I} P_i$
by the rule $(i', p') \geq (i, p)$ if and only if
$i' \geq i$ and the map $B_i \to B_{i'}$ maps $A_{i, p}$ into
$A_{i', p'}$. This is exactly the condition needed to
define a morphism $S_{i', p'} \to S_{i, p}$: namely make a commutative
diagram as above using the transition morphisms $V_{i'} \to V_i$
and $W_{i'} \to W_i$ and
the morphism $\Spec(A_{i', p'}) \to \Spec(A_{i, p})$ induced
by the ring map $A_{i, p} \to A_{i', p'}$. The relevant commutativities
have been built into the constructions.
We claim that $S$ is the directed limit of the schemes $S_{i, p}$.
Since by construction the schemes $V_i$ have limit $V$ this boils
down to the fact that $B$ is the limit of the rings $A_{i, p}$
which is true by construction. The map $\alpha : J \to I$ is given
by the rule $j = (i, p) \mapsto i$. The open subscheme $V'_j$ is
just the image of $V_i \to S_{i, p}$ above. The commutativity of
the diagrams in (5) is clear from the construction.
This finishes the proof of the lemma.
\end{proof}
\end{document}
